\documentclass[10pt,a4paper]{amsart}
\usepackage[margin=19mm]{geometry}
\usepackage{amsmath,amssymb,mathtools}
\usepackage[scr=rsfs,cal=euler]{mathalfa}
\usepackage{xfrac}
\usepackage[unicode,hidelinks,bookmarks=false]{hyperref}
\newcommand{\ie}{i.e.\ }
\newcommand{\cf}{cf.\ }
\newcommand{\resp}{respectively\ }
\newcommand{\Subsection}{\subsection}
\providecommand{\nobreakdash}{\nobreak\hbox{-}}
\setlength{\emergencystretch}{2em}
\multlinegap0pt
\usepackage{enumerate}
\usepackage[all]{xy}
\usepackage{amssymb}
\usepackage{bm}
\usepackage{tikz}
\usepackage{float}
\usepackage[normalem]{ulem}
\usepackage{mathtools}
\newtheorem{thm}{Theorem}
\newtheorem{cor}{Corollary}
\newcommand{\CF}{{\mathcal F}}
\newcommand{\RR}{{\mathbb R}}
\newcommand{\Spec}{{\rm Spec}}
\newcommand{\Star}{{\rm Star}}
\newcommand{\cC}{{\mathcal{C}}}
\newcommand{\com}[1]{{\color{brown}$\textsuperscript{+}$\marginnote{\small #1}}}
\newcommand{\trop}{\mathrm{trop}}
\title{Balancing properties of tropical moduli maps}
\author{Karl Christ, Xiang He, Ilya Tyomkin}
\date{Source: arXiv:2403.15686v2 (CC BY 4.0)}
\begin{document}
\maketitle

\noindent\emph{Source and licence.} Balancing properties of tropical moduli maps. Version arXiv:2403.15686v2 (CC BY 4.0), distributed by arXiv under the Creative Commons Attribution 4.0 International licence, http://creativecommons.org/licenses/by/4.0/, as recorded in the arXiv licence metadata for this exact version and shown on the abstract page https://arxiv.org/abs/2403.15686v2. Full licence notice: \texttt{licenses/arxiv-2403.15686-CC-BY-4.0.txt}.\par\smallskip

\noindent\textit{Editorial scope.} Counted unit: the article's cor cor:main (source0.tex lines 471--480). The article's complete original proof of it is delivered with it (source0.tex lines 482--503, one \texttt{proof} environment of 22 lines). No mathematics is written, repaired or re-derived: the statement and the proof are the article's own lines, in the article's own notation, and no retained line is substituted or re-flowed.  Retained uncounted as the source-local context the proof needs: lines 451--458.  Nothing was omitted from any retained span, so the package carries no omission record.  The retained lines cite 2 works, whose entries are transcribed verbatim from the article's own bibliography source in the e-print and delivered with short editorial labels recorded in the item's reference mapping.  No retained line contains a figure environment, an image-inclusion command or drawing code, so no image is delivered and the package declares no asset dependency.  Editorial formatting only, in addition: an AMS \texttt{amsart} preamble that restates the article's own declarations and macros used by the retained lines, namely the environments \texttt{thm}, \texttt{cor} on separate counters, together with the standard packages those lines need.  The retained environments print in this document as Theorem 1, Corollary 1 in order of appearance, whereas the article prints the same items with the numbering of its own full document.  The complete native source of the article as distributed in the arXiv e-print for this exact version is bundled unchanged as \texttt{sources/156517/source0.tex}.
\begin{thm}\label{thm:main thm}
Let $(B^0,H^0)$ be a strictly semistable pair over $K^0$ and $(\cC^0\rightarrow B^0,\sigma^0_\bullet)$ a split family of stable marked curves over $B^0$. Let $B$, $H$ and $\cC$ be the generic fibers of $B^0$, $H^0$, and $\cC^0$, respectively. Suppose $\cC^0$ is obtained as a pullback of the universal family over $\overline M$ and is smooth over $B':=B\backslash H$. Suppose furthermore that $f \colon \cC|_{B'} \rightarrow X$ is a family of parameterized curves over $B'$. Then there exists a family of parameterized tropical curves $h\: \Gamma_\Lambda\to N_\RR$ over $\Lambda:=S(B^0,H^0)$ such that 
\begin{enumerate}
\item For any $\eta \in B'(K)$, the fiber of $(\Gamma_\Lambda, h)$ over $\trop(\eta)$ is the tropicalization of $f\:\cC_\eta\to X$;% with respect to the model $\cC^0|_{\overline \eta}$\com{if it is the stable model then why to mention w.r.t.which model?}; and
\item For any face $W\in\CF(\Lambda)$ of codimension one (resp. of codimension at least two) such that the tropical curves over the interior of $W$ are weightless and $3$-valent except for at most one $4$-valent vertex, 
the induced map $\alpha\:\Lambda\to M^{\trop}_{g,n,\nabla}$ is either harmonic (resp. quasi-harmonic) or locally combinatorially surjective along $W$.
\end{enumerate}
\end{thm}

\begin{cor} \label{cor:main}
Let $X$ be a toric variety and $f\colon \cC \rightarrow X$ a family of parameterized curves defined over the field $K$. Assume that the base $B'$ of the family is quasi-projective, and consider the set of tropicalizations
$$\Sigma:=\left\{\left[ \trop(f_\eta, C_\eta)\right]\; :\; \eta\in B'(K) \right\}\subseteq M^{\trop}_{g,n,\nabla}.$$
Let $M_{[\Theta]}\subset M^{\trop}_{g,n,\nabla}$ be a stratum. Then,
\begin{enumerate}
\item The closure $\overline\Sigma\subseteq M^{\trop}_{g,n,\nabla}$ is the image of a rational polyhedral complex $\Lambda$ under a piecewise integral affine map;
\item If $M_{[\Theta]}$ is weightless and $3$-valent and $\dim(\Sigma\cap M_{[\Theta]})=\dim M_{[\Theta]}$, then $\overline M_{[\Theta]}\subseteq \overline{\Sigma}$;
\item If $M_{[\Theta]}$ is weightless and almost $3$-valent and $\overline{\Sigma}$ contains one of its adjoint weightless and $3$-valent strata, then $\overline{\Sigma}$ contains all other such strata as well. 
\end{enumerate}
\end{cor}

\begin{proof}
After replacing $B'$ with a finite covering, we may assume that the family of marked curves over $B'$ is a pullback of the universal family over the projective scheme $\overline M$. Plainly, it suffices to prove the corollary for each irreducible component of $B'$. Furthermore, we may replace $B'$ with $B'_{\rm red}$, and hence we will assume that $B'$ is integral, i.e., a quasi-projective variety. Thus, there is a projective variety $B$ containing $B'$ as a dense open subset. After replacing $B$ with the closure of the graph $B'\to \overline M \times B$, we may assume that the family of marked curves over $B'$ extends to $B$, and the latter family is also a pullback of the universal family over $\overline M$. Set $H:=B\setminus B'$. 

By \cite[\href{https://stacks.math.columbia.edu/tag/0C2S}{Tag~0C2S}]{stacks-project}, $B$ admits a flat projective model $B^0$ over $K^0$. Let $\overline M^0$ be the trivial model of $\overline M$ over $K^0$. After replacing $B^0$ with the schematic image of the graph $B\to \overline M^0 \times_{\Spec(K^0)} B^0$, which is also flat over $K^0$ by {\em loc. cit.}, we may assume that the family of marked curves extends to $B^0$. By \cite[Lemma~5.3]{dJ96}, there exist a projective alteration $\varphi_1 \colon B_1^0 \to B^0$ such that the pullback of $\cC$ to $B_1^0$ is split over $B_1^0$. Replacing $B_1^0$ with an irreducible component that dominates $B^0$ if necessary, we may assume that $B_1^0$ is integral. Since the loci $\Sigma$ associated to $B'$ and to $\varphi_1^{-1}(B')$ coincide, we may assume that $B_1^0=B^0$ and hence the family $\cC$ is split over $B^0$.
Finally, by \cite[Theorem~6.5]{dJ96}, there exists an alteration $\varphi_2 \colon B_2^0 \to B^0$ such that $B_2^0$ is projective over $K^0$, the generic fiber of $B_2^0\to\Spec(K^0)$ is irreducible, and the pair $\left(B_2^0, \overline{\varphi_2^{-1}(H)}_{\rm red}\right)$ is strictly semistable. Here $\overline{\varphi_2^{-1}(H)}_{\rm red}$ is the closure of $\varphi_2^{-1}(H)_{\rm red}$ in $B_2^0$.
As before, we may assume that $B_2^0=B^0$. 

Denote by $H^0$ the union of the horizontal components of $H\cup \tilde{B}$. Then $(B^0, H^0)$ is a strictly semistable pair, and the family of marked curves extends to $B^0$ by pulling back the universal family over $\overline M$. Therefore, we may assume that we are given all the data $B^0, H^0, \cC^0, \sigma^0_\bullet, B'$ satisfying the assumptions of Theorem~\ref{thm:main thm}. Consider the tropicalization $h\: \Gamma_\Lambda\to N_\RR$ and the induced map $\alpha\:\Lambda\to M^{\trop}_{g,n,\nabla}$ provided by Theorem~\ref{thm:main thm}.

(1) Since $\alpha$ is piecewise integral affine, it follows that $\alpha(\Lambda)\subseteq M^{\trop}_{g,n,\nabla}$ is closed. Now, by Theorem~\ref{thm:main thm} (1), $\Sigma\subset \alpha(\Lambda)$, and therefore $\overline\Sigma\subset \alpha(\Lambda)$. Furthermore, $\alpha^{-1}(\Sigma)$ is dense in $\Lambda$. Thus, $\Sigma$ is dense in $\alpha(\Lambda)$, and hence $\overline\Sigma=\alpha(\Lambda)$ as asserted.

(2) Set $k:=\dim M_{[\Theta]}$, and assume to the contrary that $\overline M_{[\Theta]}\nsubseteq \alpha(\Lambda)$. By (1), $\alpha(\Lambda)\subseteq M^{\trop}_{g,n,\nabla}$ is closed, and therefore $M_{[\Theta]}\nsubseteq \alpha(\Lambda)$. Furthermore, there is a face $W\in \CF(\Lambda)$ such that $\alpha(W^\circ)$ has dimension $k-1$, $\alpha\left(\Star(W^\circ)\right)$ has dimension $k$, and $\alpha(W^\circ)$ belongs to the relative boundary of $\alpha(\Lambda)$ in $M_{[\Theta]}$.

After replacing $W$ with one of its faces, we may assume that $\dim(W)=k-1$. Indeed, since $\dim(\alpha(W^\circ))=k-1$, the dimension of $W$ is at least $k-1$. If it is greater than $k-1$, then the kernel of the differential $d\alpha$ along $W$ is non-trivial. But $\Lambda$ is the skeleton of a strictly semistable pair, and therefore $W$ contains no lines. Thus, $\alpha(W)=\alpha(\partial W)$, and by induction, we can find a face of $W$ of dimension $k-1$ satisfying the requirements.

By construction, $\alpha\left(\Star(W^\circ)\right)$ belongs to a half-space supported by $\alpha(W^\circ)$ and the differential $d\alpha$ is not identically zero on $\Star(W^\circ)$. Therefore $\alpha$ is not quasi-harmonic along $W$ contradicting Theorem~\ref{thm:main thm} (2).

(3) Let $M_{[\Theta']}$ be a weightless and $3$-valent stratum adjacent to $M_{[\Theta]}$ that is contained in $\overline\Sigma$, and $M_{[\Theta'']}$ be another weightless and $3$-valent stratum adjacent to $M_{[\Theta]}$. Then there exists a face $W'\in\mathcal F(\Lambda)$ and a face $W$ of $W'$ of codimension one such that $\alpha\left((W')^\circ\right)\subseteq M_{[\Theta']}$ has dimension $\dim M_{[\Theta']}$, and $\alpha(W^\circ)\subseteq M_{[\Theta]}$. Thus, $\alpha$ is locally combinatorially surjective along $W$ by Theorem~\ref{thm:main thm} (3). In particular, there exists a face $W''\in\mathcal F(\Lambda)$ containing $W$ as a codimension-one face such that $\alpha\left((W'')^\circ\right)\subseteq M_{[\Theta'']}$, and hence 
$$\dim \alpha\left(W''\right) \ge \dim \alpha(W)+1=\dim\left(M_{[\Theta'']}\right)$$
since $\alpha|_{W''}$ is an affine map and $\alpha(W)\subsetneq\alpha(W'')$.
Therefore, $\overline M_{[\Theta'']}\subseteq \overline{\Sigma}$ by (2), and we are done.
\end{proof}

\begin{thebibliography}{99}
\bibitem[dJ96]{dJ96} A.~J. de~Jong, \emph{Smoothness, semi-stability and alterations}, Inst. Hautes \'{E}tudes Sci. Publ. Math. (1996), no.~83, 51--93. \MR{1423020}
\bibitem[stacks]{stacks-project} The {Stacks project authors}, \emph{The stacks project}, \url{https://stacks.math.columbia.edu}, 2025.
\end{thebibliography}
\end{document}
